
\documentclass[11pt]{article}
\usepackage[margin=1in]{geometry}
\usepackage{amsmath, amssymb}
\usepackage{graphicx}
\usepackage{booktabs}
\usepackage[hidelinks]{hyperref}

\title{Robust Benchmarking via Multi-Sample Evaluation}
\author{(Your Name)}
\date{\today}

\begin{document}
\maketitle

% This file is self-contained and should compile with pdflatex/xelatex.
% It does not assume a specific paper organization.
% You can keep everything under this single section or move parts elsewhere later.

\section*{Introduction}

Modern model evaluation often relies on a single held-out test set and a point estimate such as accuracy or F1. This practice can be misleading for two reasons. First, one-shot testing confounds the true performance of a model with the idiosyncrasies of a particular sample. Second, extremely high scores (e.g., 99--100\%) are frequently artifacts of unrepresentative class balance or leakage rather than genuine generalization. To obtain a more faithful picture, we argue that a benchmark should be constructed with \emph{explicit data composition} and \emph{multi-sample testing} at its core, and that reporting should center on \emph{ranges} and \emph{uncertainty} rather than single numbers.

\paragraph{Data composition and realism.}
Let $\mathcal{D}$ denote a pool of labeled examples partitioned into positives $\mathcal{D}^{+}$ and negatives $\mathcal{D}^{-}$.
A realistic benchmark specifies a target class ratio $(\pi^{+}, \pi^{-})$ with $\pi^{+} + \pi^{-} = 1$ (e.g., $\pi^{+}=0.8,\ \pi^{-}=0.2$ or $\pi^{+}=0.9,\ \pi^{-}=0.1$), reflecting the expected operational prevalence.
Benchmarks may be constructed per dataset (single-source baseline) or by mixing multiple sources (pooled baseline) to capture heterogeneity.

\paragraph{Repeated out-of-sample evaluation.}
Given a model $f$ and a metric $m(\cdot)$ (e.g., Accuracy, Precision, Recall, F1, AUROC), we draw $R$ independent \emph{out-of-sample} test sets $\{\mathcal{T}_r\}_{r=1}^{R}$ from $\mathcal{D}$ according to $(\pi^{+}, \pi^{-})$, each of size $n$.
We compute $m_r = m(f;\mathcal{T}_r)$ for every replicate.
The reporting target becomes the sampling distribution summary:
\begin{equation*}
\bar{m}=\frac{1}{R}\sum_{r=1}^{R} m_r,\qquad
s_m=\sqrt{\frac{1}{R-1}\sum_{r=1}^{R} (m_r - \bar{m})^2}.
\end{equation*}
Instead of a single point, we present $\bar{m} \pm s_m$ and a $(1-\alpha)$ confidence interval (CI).
For large $R$, a normal-approximate CI is
\begin{equation*}
\text{CI}_{1-\alpha}: \quad \bar{m} \pm z_{1-\alpha/2}\cdot \frac{s_m}{\sqrt{R}},
\end{equation*}
or the Student-$t$ variant when $R$ is modest.
This reframes evaluation as \emph{performance with quantified uncertainty}.

\paragraph{Multiple baselines and factorization of variability.}
Beyond a single composition, we consider a grid of benchmark settings:
(i) \textbf{per-dataset} baselines (each source independently),
(ii) \textbf{mixed} baselines (sources pooled at controlled proportions), and
(iii) \textbf{composition sweeps} across $\pi^{+}\in\{0.5,0.8,0.9\}$ (and $\pi^{-}=1-\pi^{+}$).
For each setting, we repeat the sampling protocol to estimate $(\bar{m}, s_m, \text{CI})$.
This exposes sensitivity to dataset origin and class balance, two dominant axes of real-world variation.

\paragraph{Comparing multiple models.}
Let $\mathcal{F}=\{f^{(1)},\dots,f^{(K)}\}$ denote competing models (e.g., distinct LMs or systems).
We evaluate each $f^{(k)}$ under the same sampling plan and report the collection $\{(\bar{m}^{(k)}, s_m^{(k)}, \text{CI}^{(k)})\}_{k=1}^{K}$ per benchmark setting.
Pairwise differences can be tested by resampling-aware procedures.
For example, if $m_r^{(a)}$ and $m_r^{(b)}$ are paired by identical $\mathcal{T}_r$, one may use a paired $t$-test on $\Delta_r = m_r^{(a)} - m_r^{(b)}$, reporting
\begin{equation*}
\bar{\Delta}=\frac{1}{R}\sum_{r=1}^{R} \Delta_r,\qquad
s_\Delta=\sqrt{\frac{1}{R-1}\sum_{r=1}^{R}(\Delta_r-\bar{\Delta})^2},
\end{equation*}
and a CI for $\bar{\Delta}$.
When distributional assumptions are doubtful, a nonparametric alternative (e.g., permutation or bootstrap) provides robust $p$-values and CIs.

\paragraph{From single numbers to ranges and decisions.}
Operational decisions rarely hinge on a single-point score; they trade off risk and benefit under uncertainty.
Therefore, we encourage reporting \emph{performance bands} (e.g., violin/interval plots over replicates), together with the estimated probability that one model outperforms another by a practically meaningful margin (e.g., $\Pr[\bar{\Delta}>\delta]$ for a chosen effect size $\delta$).
This connects statistical significance to \emph{practical significance}.

\paragraph{Controlling for synthetic and ``peaky'' cases.}
Where positive or negative examples can be curated or synthesized, we incorporate stress sets that are intentionally ``hard'' (e.g., distributional shifts, class-ambiguous or high-confusability items).
These stress sets are included as additional benchmark cells, evaluated with the same replicate protocol, to surface brittleness that average-case metrics may hide.

\paragraph{Recommended reporting template (drop-in).}
For each benchmark cell (dataset source, class ratio, stress/regular), report:
\begin{enumerate}
  \item Replicate count $R$ and test size $n$; sampling recipe (with/without replacement).
  \item Class composition $(\pi^{+},\pi^{-})$ and any mixing proportions across sources.
  \item Metric(s) $m$ with $\bar{m}$, $s_m$, and $(1-\alpha)$ CI.
  \item Model comparisons: $\bar{\Delta}$, CI, and $p$-value (paired if applicable).
  \item Visual summary: interval/violin plots across replicates; optionally calibration or PR/ROC curves aggregated over replicates.
\end{enumerate}

\paragraph{Notation cheat sheet (for quick reuse).}
\begin{itemize}
  \item $\mathcal{D}$: labeled pool; $\mathcal{D}^{+}$ positives, $\mathcal{D}^{-}$ negatives.
  \item $(\pi^{+},\pi^{-})$: target class ratio; $n$: test size; $R$: replicate count.
  \item $\mathcal{T}_r$: $r$-th out-of-sample test set; $m_r$: metric on $\mathcal{T}_r$.
  \item $\bar{m}, s_m$: replicate mean and standard deviation; $\text{CI}$: confidence interval.
  \item $f^{(k)}$: $k$-th model; $\Delta_r$: replicate-wise paired difference between two models.
\end{itemize}

\paragraph{Why this matters.}
This protocol converts evaluation from a one-time number into a distributional characterization that respects data realism (class balance), assesses stability (replicates), and supports principled comparison (CIs, $p$-values, and effect sizes). It reduces the risk of overclaiming from flattering splits and provides actionable uncertainty that downstream stakeholders can plan against.

\paragraph{Question and Discussion}
see benchmark\_discussion.pdf
\end{document}

